% Medication entries centralized from the alphabetical chapters.
\section*{Generic Medicines and Medication Terms}
\medterm{5-Alpha Reductase Inhibitor} A medicine that lowers production of dihydrotestosterone, a hormone that affects the prostate and hair follicles. It is used mainly for an enlarged prostate and selected types of hair loss.

\medterm{Abortifacient} A medicine or substance used to end a pregnancy. Some block progesterone, others cause uterine contractions and cervical opening, and certain medicines stop rapidly dividing pregnancy tissue in selected ectopic pregnancies.

\medterm{Abortion - Medication} Medication abortion uses prescribed medicines, commonly mifepristone followed by misoprostol or misoprostol alone, to end an early intrauterine pregnancy; confirmation of completion and assessment for ectopic pregnancy are important.

\medterm{Abuse Steroid (Anabolic Steroid Abuse)} Alternate index label for Anabolic Steroid Abuse.

\medterm{Accident And Emergency Medicine (A\&E Medicine)} The medical specialty concerned with the immediate assessment, resuscitation, diagnosis, and initial treatment of acute illness and injury.

\synonyms
A\&E Medicine

\medterm{ACE Inhibitor} A medicine that relaxes blood vessels and reduces strain on the heart by lowering production of angiotensin II. It treats high blood pressure, heart failure, and some kidney disease; cough, high potassium, kidney problems, or swelling can occur.

\medterm{Acitretin} An oral retinoid used mainly for severe psoriasis and selected keratinization disorders. It is highly teratogenic and requires strict pregnancy precautions and monitoring for adverse effects.

\medterm{Adenovirus Vaccine} A vaccine used primarily to prevent adenovirus types 4 and 7 among selected
military personnel. It is an oral, live vaccine and is not part of routine vaccination for the general
public. Vaccine use and eligibility depend on current public-health recommendations.

\medterm{Adverse Drug Event} Harm related to the use of a medication, including harm caused by an adverse drug reaction, medication error,
overdose, or inappropriate use.

\medterm{Adverse Drug Reaction} An unwanted and harmful response to a medicine taken at a usual dose for prevention, testing, or treatment. It may result from the medicine’s expected effects, an allergy, an interaction, or an unusual individual response.

\medterm{Agent, Tocolytic} A medicine that reduces uterine contractions, used selectively to delay preterm birth long enough for corticosteroids, transfer, or other preparations to take effect.

\medterm{Agonist-Antagonist} A medicine that activates some receptors while blocking others, or activates a receptor only partly. Its mixed action can produce a desired effect while limiting another effect; examples include some opioid medicines.

\medterm{AIDS Vaccine} A vaccine designed to prevent HIV infection or reduce its effects; no generally effective licensed vaccine currently exists.

\medterm{Alendronate} A bisphosphonate that reduces bone resorption and is used to treat or prevent osteoporosis and selected other disorders of excessive bone turnover.

\synonyms
Alendronic acid; oral bisphosphonate therapy.




\medterm{Anabolic Steroid Abuse} Anabolic-androgenic steroid misuse can cause infertility, testicular shrinkage, acne, mood changes, high blood pressure, adverse cholesterol changes, liver injury, and cardiovascular events. Stopping may produce withdrawal symptoms and is safest with medical support.

\synonyms
Abuse Steroid (Anabolic Steroid Abuse)

\medterm{Analgesic} A medicine that relieves pain without necessarily causing loss of consciousness. Analgesics include nonopioid medicines, opioids, and other medicines chosen according to the type and severity of pain.

\medterm{Analgesic Nephropathy} Chronic tubulointerstitial renal injury associated with prolonged or excessive use of
analgesic medicines. It may cause impaired urinary concentration, sterile pyuria,
haematuria, papillary necrosis, hypertension, and progressive chronic kidney disease.

\medterm{Analgesics} Medicines used to relieve pain. They include nonopioid analgesics, opioid analgesics, and adjuvant medicines used for selected pain conditions.

\synonyms
Painkillers

\medterm{Anesthetic} A medicine or technique that numbs part or all of the body. Depending on the method, a person may remain awake or become unconscious during a procedure; breathing, blood pressure, and other vital functions may need monitoring.

\medterm{Angiogenesis Inhibitor} A medicine that slows the growth of new blood vessels. It may be used in selected cancers or eye diseases, where abnormal blood-vessel growth helps the disease progress; bleeding, high blood pressure, and other effects may occur.

\synonyms
Antiangiogenic agent

\medterm{Angiotensin Converting-Enzyme Inhibitor} A medicine that lowers the production of angiotensin II, a hormone that narrows blood vessels and promotes salt retention. It can treat high blood pressure, heart failure, and some kidney disease but may cause cough, high potassium, kidney problems, or swelling.

\medterm{Angiotensin Receptor Blocker} A medicine that blocks angiotensin II at the AT1 receptor, reducing vasoconstriction and aldosterone effects; it is used for hypertension, heart failure, and selected kidney disease.

\medterm{Anti-Infective} A medicine or other agent used to prevent or treat an infection. The choice depends on whether the cause is a bacterium, virus, fungus, or parasite, as well as the infection site and the person’s health.

\medterm{Anti-Inflammatory} A medicine or treatment that reduces inflammation, which is the body’s response to injury or illness. Different types work in different ways and may have important effects on the stomach, kidneys, heart, or immune system.

\medterm{Antiarrhythmic} A medicine used to prevent or treat an abnormal heart rhythm. Different medicines alter the heart’s electrical signals in different ways, and the choice depends on the rhythm, heart condition, kidney function, and other risks.

\medterm{Antibiotic} A medicine that kills bacteria or stops them from multiplying. Antibiotics do not treat viral infections and should be selected for the suspected bacteria, infection site, allergies, and resistance risk.

\medterm{Antibiotic Cycling} A strategy in which the preferred antibiotic for a particular clinical use is changed at
planned intervals. It is intended to limit selection of resistant organisms, although its
benefit depends on local conditions and implementation.

\medterm{Antibiotic De-Escalation} The process of narrowing antimicrobial spectrum based on microbiological results and
clinical response. De-escalation aims to reduce selective pressure for resistance while
maintaining therapeutic efficacy.

\medterm{Antibiotic Mixing} A strategy in which different antibiotics are assigned concurrently to different patients
with the same indication. It is intended to reduce uniform selection pressure for
resistance in a patient population.

\medterm{Antibiotic Resistance} The ability of bacteria to survive or grow despite an antibiotic that would normally kill them or stop their growth. It can result from genetic changes and the spread of resistance genes; unnecessary or incorrect antibiotic use can accelerate it.

\synonyms
Resistance Antibiotic (Antibiotic Resistance)

\medterm{Antibiotic Resistance Genes} Genes that help bacteria survive antibiotics. They may change the antibiotic’s target, destroy or modify the medicine, or remove it from the cell. Bacteria can inherit these genes or acquire them from other bacteria.

\medterm{Antibiotic Stewardship} Coordinated efforts to use antibiotics appropriately by selecting the right medicine, dose, route, and duration while reducing unnecessary treatment and resistance.

\medterm{Antibiotic-Associated Colitis} Alternate Mayo Clinic label; see \textit{Pseudomembranous colitis}.

\medterm{Antibiotic-Associated Diarrhea} Loose or frequent stools that occur during or soon after antibiotic treatment.

\synonyms
Diarrhea, Antibiotic-Associated

\medterm{Anticholinergic Cognitive Burden Scale} A medication-review tool that estimates the total anticholinergic load of a patient's
medicines. It supports risk assessment and deprescribing but does not replace
individualized clinical judgment.

\medterm{Anticoagulant} A medicine that slows the body’s blood-clotting process. It helps prevent or treat clots in veins, lungs, heart, or arteries but can cause bruising or serious bleeding.

\medterm{Anticoagulant Agent} A drug that reduces blood clot formation by interfering with the coagulation cascade; examples include heparin, warfarin, and direct oral anticoagulants.

\medterm{Anticoagulant Monitoring} Laboratory assessment used to measure the effect or safety of an anticoagulant. The
appropriate test depends on the drug and may include aPTT, anti-Xa activity, INR, or
measurements of kidney function and blood counts.

\medterm{Anticoagulant Rodenticides Poisoning} Poisoning from rat or mouse bait that interferes with blood clotting. Symptoms such as easy bruising, nosebleeds, blood in urine or stool, or weakness may appear later; suspected exposure requires immediate poison-control or emergency advice.

\medterm{Anticonvulsant} A medicine used to prevent or control seizures. Some are also used for nerve pain, migraine prevention, or mood disorders; choice depends on the condition, other medicines, pregnancy, and side effects.

\synonyms
Antiseizure medicine

\medterm{Anticonvulsant Mood Stabilizers} Antiepileptic drugs used as mood stabilizers in bipolar disorder. Valproate (divalproex)
is effective for acute mania and rapid cycling. Carbamazepine is effective for acute
mania but has numerous drug interactions via CYP3A4 induction. Lamotrigine is effective
for bipolar depression prevention but not acute mania.

\medterm{Antidepressant} A medicine used to treat depressive disorders and some other conditions involving mood, anxiety, pain, or sleep.

\medterm{Antidiabetic Agent} A medicine that lowers blood glucose or improves glucose regulation in diabetes, such as insulin, metformin, sulfonylureas, or GLP-1 receptor agonists.

\medterm{Antidiarrheal} A medicine that reduces the frequency or amount of diarrhea. It may be unsafe when stools contain blood, there is high fever or severe infection, or the bowel is blocked, so the cause should be considered first.

\medterm{Antidiarrheal Medicine Overdose} Toxic effects caused by excessive exposure to Antidiarrheal medicine. Severity depends on the dose, route, timing, coexisting conditions, and clinical findings; suspected poisoning requires prompt assessment.

\medterm{Antiemetic} A medicine that prevents or reduces nausea and vomiting. The choice depends on the cause, such as infection, migraine, pregnancy, surgery, or cancer treatment, and some medicines can cause sleepiness or movement problems.

\medterm{Antiestrogen} A drug that blocks estrogen effects or reduces estrogen production; it may be used in hormone-sensitive breast cancer or to alter reproductive hormone signaling.

\medterm{Antifungal} A medicine that kills fungi or stops them from growing. It may be applied to the skin or given by mouth or injection; the choice depends on the fungus, the body site, and the severity of infection.

\medterm{Antifungal Agent} A drug or chemical that kills fungi or inhibits their growth, used to treat fungal infections of the skin, mucosa, or internal organs.

\medterm{Antifungal Drug Classes} Major classes of antifungal agents include polyenes (amphotericin B), azoles
(fluconazole, voriconazole), echinocandins (caspofungin, micafungin), and allylamines
(terbinafine). Polyenes bind to ergosterol in fungal cell membranes. Azoles inhibit
ergosterol synthesis. Echinocandins inhibit cell wall synthesis.

\medterm{Antifungal Medication} A medicine used to treat fungal infection; the choice and route depend on the organism and site, and examples include azoles, terbinafine, polyenes, and echinocandins.

\medterm{Antihistamine} A medicine that blocks histamine, a chemical involved in allergies. It can relieve sneezing, itching, hives, or watery eyes; older types may cause sleepiness, dry mouth, blurred vision, urinary retention, or confusion.

\medterm{Antihypertensive} A medicine used to lower blood pressure. Classes include ACE inhibitors, angiotensin-receptor
blockers, calcium-channel blockers, thiazide-type diuretics, beta-blockers, and other
agents. Selection depends on blood pressure, kidney function, cardiovascular disease,
pregnancy, and other clinical factors.

\medterm{Antihypertensive Agent} A medicine or treatment that lowers high blood pressure. The choice depends on age, kidney function, heart disease, pregnancy, other medicines, and how high the pressure is.

\synonyms
Blood-pressure medicine

\medterm{Antimalarial} A drug used to prevent or treat malaria caused by Plasmodium parasites; selection depends on species, disease severity, resistance, and geographic exposure.

\medterm{Antimicrobial Drug} A medicine used to treat infection by bacteria, viruses, fungi, or parasites; the drug is selected according to the organism, infection site, and patient factors.

\medterm{Antimicrobial Drug Interactions} Changes in the effect or level of an infection medicine caused by another medicine, food, or supplement. Some combinations reduce treatment effectiveness or increase toxicity, so the full medicine list should be checked.

\medterm{Antineoplastic} A medicine or treatment that kills cancer cells or slows their growth. It may be used alone or with surgery, radiation, immunotherapy, or other treatment and can also affect healthy fast-growing cells.

\medterm{Antiparasitic} A medicine used to prevent or treat infection caused by parasites. The choice depends on the parasite, body site, disease severity, resistance patterns, pregnancy status, and other patient factors.

\medterm{Antiparkinsonian} A medicine used to improve movement or other symptoms of Parkinson disease. Medicines may replace dopamine, imitate its effects, or alter other brain signals; benefits and adverse effects vary with the drug and the person.

\medterm{Antiplatelet Agent} A medicine that makes blood platelets less likely to clump together. It lowers the risk of clots in arteries but can increase bruising and serious bleeding.

\synonyms
Antiplatelet medicine

\medterm{Antiprotozoal Drug} A medicine used to prevent or treat disease caused by protozoa, including drugs directed against malaria, intestinal protozoa, or tissue parasites.

\medterm{Antipruritic} A medicine or skin treatment that relieves itching. It may work by reducing inflammation, blocking allergy signals, numbing the skin, or changing nerve signals; treatment should also address the cause of the itching.

\medterm{Antipsychotic} A medicine used to reduce hallucinations, delusions, severe agitation, or mania and to treat some mood disorders. Possible effects include sleepiness, movement problems, weight or blood-sugar changes, and rarely a dangerous fever-and-stiffness reaction.

\medterm{Antipsychotic Agent} A medicine used to reduce hallucinations, fixed false beliefs, severe agitation, or disorganized thinking. Some are also used for bipolar disorder or selected behavioral symptoms and can cause sleepiness, movement problems, or metabolic effects.

\synonyms
Antipsychotic

\medterm{Antipsychotic-Induced Parkinsonism} Drug-induced bradykinesia, rigidity, tremor, or postural instability caused by
dopamine-receptor blockade, usually from antipsychotic treatment. Management involves
reviewing the medication, dose, alternatives, and symptomatic treatment when
appropriate.

\medterm{Antipyretic} A medicine that lowers fever by reducing the brain’s temperature-setting response. It may also relieve pain, but it does not treat every cause of fever; the cause and warning symptoms still need attention.

\medterm{Antispasmodic} A medicine that relaxes involuntary muscle and may reduce cramping in the intestine, bladder, or other organs. Some types can cause dry mouth, blurred vision, constipation, urinary retention, or confusion.

\synonyms
Antispasmodics

\medterm{Antithyroid} A medicine that lowers thyroid-hormone production or release, or reduces the effects of thyroid hormone. Treatment is selected according to the cause of hyperthyroidism, disease severity, pregnancy status, and other clinical factors.

\medterm{Antithyroid Drug} A medicine that reduces thyroid hormone production or release, used to treat hyperthyroidism; examples include methimazole, propylthiouracil, and radioactive iodine.

\medterm{Antitussive} A medicine that reduces the urge or reflex to cough. It may help a troublesome dry cough, but suppressing a productive cough can be unhelpful and serious causes of cough still need assessment.

\medterm{Antiviral} A medicine that slows or stops viruses from multiplying. It works only against certain viruses and is chosen according to the infection, timing, resistance, and the patient’s health.

\medterm{Antiviral Drug Classes} Groups of antiviral medicines classified by the step of viral growth they block. Some stop copying of viral genetic material, some block viral proteins or entry into cells, and others prevent newly formed viruses from being released.

\medterm{Anxiolytic} A medicine or other treatment that reduces anxiety. Choice depends on the anxiety disorder, urgency, coexisting
conditions, adverse effects, and patient preference; psychological therapies and, when
appropriate, antidepressant medicines are commonly used for persistent anxiety.

\medterm{ARB-Neprilysin Inhibitor} Sacubitril/valsartan combines an ARB with a neprilysin inhibitor, increasing natriuretic
peptide levels while blocking the renin-angiotensin system. The PARADIGM-HF trial
demonstrated superior reduction in cardiovascular death and heart failure
hospitalization compared to enalapril.


\medterm{Aromatase Inhibitor} A class of medications that inhibit aromatase, the enzyme that converts androgens
(androstenedione, testosterone) to estrogens (estrone, estradiol) in peripheral tissues
(adipose, breast, muscle, bone). Used primarily in hormone receptor-positive breast
cancer in postmenopausal women.

\medterm{Aspirin} Acetylsalicylic acid, a salicylate that relieves pain and fever, reduces inflammation, and irreversibly inhibits platelet aggregation; it can cause gastrointestinal bleeding and is generally avoided in children with viral illness.

\medterm{Asthma OTC Medication (Asthma: Over The Counter Treatment)} Alternate index label for Asthma: Over The Counter Treatment.

\medterm{Attenuated Vaccine} A vaccine made with a weakened microorganism that retains the ability to stimulate immunity.

\synonyms
Live attenuated vaccine

\medterm{Bacterial Vaccine} A vaccine made from killed bacteria, weakened bacteria, purified bacterial parts, or bacterial toxins. It trains the immune system to prevent specific infections; protection and dosing depend on the vaccine and age group.

\medterm{BCG Vaccine} A vaccine made from a weakened relative of the tuberculosis bacterium. It helps protect children from severe tuberculosis, especially meningitis, but protection against lung tuberculosis varies; use depends on local risk and policy.

\medterm{Behavioral Medicine} An interdisciplinary field integrating behavioral science with medical practice to prevent, diagnose, and treat illness. It addresses psychological factors such as stress, coping, adherence, and health behaviors that influence disease outcomes.

\medterm{Benzodiazepines} Medicines that strengthen the brain’s main calming signal. They can reduce anxiety, stop some seizures, relax muscles, or cause sleepiness, but can also cause dependence, falls, confusion, and dangerously slow breathing, especially with opioids or alcohol.


\medterm{Beta Agonists} Drugs that activate beta-adrenergic receptors, producing effects such as relaxation of airway smooth muscle or increased cardiac stimulation depending on receptor selectivity.

\synonyms
Beta-adrenergic agonists; beta-agonist drugs.

\medterm{Beta Blocker} A drug that blocks beta-adrenergic receptors, lowering sympathetic effects such as heart rate and contractility; beta blockers are used for conditions including hypertension, angina, arrhythmias, and heart failure.

\medterm{Beta Blockers} A class of medicines that block the effects of adrenaline and related hormones at beta receptors. They usually slow the heart rate, reduce the force of heart contraction, and lower blood pressure. Depending on the medicine, they may be used for high blood pressure, angina, abnormal heart rhythms, heart failure, migraine prevention, tremor, or some symptoms of anxiety.

\medterm{Beta-Adrenergic Blocker} A drug that antagonizes beta-adrenergic receptors, reducing heart rate, myocardial
contractility, and renin release. It is used in selected patients with hypertension,
angina, arrhythmias, post-myocardial infarction disease, and heart failure; abrupt
withdrawal may worsen ischemia.

\medterm{Beta-Agonist} A medicine that stimulates beta receptors. Beta-2 agonists relax airway muscles to improve breathing, while other types increase heart activity; possible effects include tremor, fast heartbeat, or low potassium.

\medterm{Beta-Blocker Overdose} Toxic exposure to a beta-adrenergic blocker that may cause slow heart rate, low blood pressure, low blood glucose, bronchospasm, seizures, or altered consciousness. It requires urgent poison-control and medical management.

\medterm{Beta-Lactamase Inhibitor Combinations} Combinations of beta-lactam antibiotics with beta-lactamase inhibitors that extend
activity against beta-lactamase producing organisms. The major combinations include
amoxicillin-clavulanate, ampicillin-sulbactam, piperacillin-tazobactam,
cefoperazone-sulbactam, and ceftolozane-tazobactam.


\medterm{C1 Esterase Inhibitor} A blood protein that helps control inflammation, swelling, and clotting reactions. Too little or abnormal activity can cause hereditary angioedema, with repeated attacks of swelling in the skin, bowel, or airway.

\medterm{Calcium Channel Blocker} A medicine that reduces calcium movement into heart and blood-vessel muscle cells. It may relax arteries, lower blood pressure, reduce chest pain, or slow some abnormal heart rhythms; ankle swelling, constipation, dizziness, or a slow heartbeat may occur.

\medterm{Calcium Channel Blocker Overdose} Toxic exposure to a calcium-channel blocker that may cause low blood pressure, slow heart rate, conduction abnormalities, high blood glucose, or shock. It requires urgent medical and poison-control management.

\medterm{Calcium Channel Blockers} Medicines that reduce the movement of calcium into heart and blood-vessel muscle cells. Dihydropyridines, such as amlodipine and nifedipine, mainly relax arteries and are used for high blood pressure and angina. Other calcium channel blockers, such as diltiazem and verapamil, also slow the heart rate and electrical conduction and may be used for certain abnormal heart rhythms.

\medterm{Camptothecin Analogue} A cancer medicine related to camptothecin that blocks an enzyme cancer cells need to copy DNA. It can damage rapidly dividing cells and commonly causes low blood counts, diarrhea, nausea, or other gastrointestinal effects.

\medterm{Cancer Vaccine} A vaccine designed either to prevent infections that can cause cancer or to help the immune system attack an existing cancer. Its use depends on the specific vaccine and cancer; it does not replace other cancer treatment when needed.

\medterm{Cardiotonic Agent} A drug or substance that increases the force or efficiency of cardiac contraction, used selectively in heart failure or shock because benefits must be balanced against arrhythmia and other risks.


\medterm{Carvedilol Phosphate} The phosphate formulation of carvedilol, a medicine that blocks beta and alpha-1 receptors. It lowers heart rate and blood pressure and is used for heart failure and hypertension; dizziness, slow heartbeat, and worsening fluid retention can occur.



\medterm{Cebocephaly} A rare birth defect in which the middle part of the face is severely underdeveloped, sometimes with a single nostril or a nose-like structure between closely spaced eyes. It usually occurs with major brain and facial-development abnormalities.

\medterm{Celecoxib} A selective COX-2 anti-inflammatory medicine used for pain and arthritis. It can cause stomach, kidney, blood-pressure, and heart problems and may increase cardiovascular risk, especially with long-term use or in people with existing risk factors.

\medterm{Checkpoint Inhibitor} A medicine that blocks an inhibitory immune-checkpoint pathway, helping selected immune cells recognize and attack tumor cells. Examples target PD-1, PD-L1, or CTLA-4.

\medterm{Cholesterol - Drug Treatment} Medicines used to lower harmful blood cholesterol or triglycerides and reduce cardiovascular risk. Statins are used most often; other options may be added when needed, according to cholesterol levels, overall risk, other illnesses, and treatment tolerance.

\medterm{Cholinesterase Inhibitor} A medicine that slows the breakdown of acetylcholine, a chemical used for nerve signalling. It can improve nerve communication in selected conditions but may cause nausea, diarrhoea, slow heartbeat, sweating, or muscle cramps.

\medterm{Cold And Cough Medicine For Infants And Children} Many nonprescription cold and cough medicines are unsafe or ineffective for young children and can cause sedation, agitation, abnormal heart rate, overdose, or breathing problems. Age-appropriate fluids, saline, humidification, and clinician-guided treatment are safer approaches.


\medterm{Conjugate Vaccine} A vaccine that joins a bacterial sugar to a carrier protein so the immune system makes a stronger, longer-lasting response. This design is especially useful in infants and young children.

\synonyms
Protein-conjugate vaccine; conjugated polysaccharide vaccine.

\medterm{Conscious Sedation} Medicine given to make a person relaxed and drowsy during a procedure while allowing them to respond to voice or touch and usually breathe independently. Breathing, blood pressure, and alertness must be monitored because deeper sedation can occur.

\medterm{Controlled Substance} A medicine or chemical whose production, supply, possession, or use is restricted by law because it may cause dependence, misuse, overdose, or other serious harm. Rules differ between countries and substances.

\medterm{COX Inhibitor} A class of medicines that reduce production of chemicals involved in pain, fever, and inflammation. They can irritate the stomach, affect kidney function, raise blood pressure, or increase bleeding and cardiovascular risk.

\medterm{Cox-1 Inhibitor} A medicine that blocks cyclooxygenase-1, an enzyme involved in pain, stomach protection, kidney function, and platelet clotting. Blocking it may relieve pain but can increase stomach irritation, bleeding, or kidney problems.

\medterm{Crack (Drug)} A smokable, rapidly acting form of cocaine; it can cause dependence, severe cardiovascular effects, seizures, stroke, psychiatric symptoms, and overdose.

\medterm{Cutaneous Drug Eruptions} Skin rashes or other skin reactions caused by medicines. They may appear as a widespread red rash, hives, sun-sensitive areas, or a lichen-planus-like eruption. Blisters, skin peeling, mouth sores, fever, or facial swelling require urgent assessment.

\medterm{Cyclooxygenase Inhibitor} A medicine that reduces production of prostaglandins and may relieve pain and inflammation.


\medterm{Designer Steroid} A chemically modified anabolic-androgenic steroid created to produce performance-enhancing effects or evade detection; products may be unapproved, contaminated, and medically dangerous.

\medterm{Diarrhea, Antibiotic-Associated} Alternate index label; see \textit{Antibiotic-associated diarrhea}.


\medterm{Digitalis} A group of heart medicines from foxglove plants, including digoxin. They can strengthen heart contraction and slow some heart rhythms but have a narrow safety range and may cause dangerous toxicity.

\medterm{Diphtheria Vaccine} A vaccine containing an inactivated form of diphtheria toxin. It trains the immune system to neutralize the toxin and is usually given in combination vaccines with protection against tetanus and pertussis.

\medterm{Disease-Modifying Antirheumatic Drug} A medicine that reduces immune-driven inflammation and helps prevent lasting joint or organ damage, rather than only relieving symptoms. Different types are chosen according to the disease, severity, other conditions, and monitoring needs.

\medterm{Disulfiram} A medicine used as an adjunct for alcohol-use disorder that inhibits aldehyde dehydrogenase and can cause an unpleasant reaction if alcohol is consumed.

\medterm{Diuretic} A medicine that helps the kidneys remove extra salt and water in urine. It can treat swelling, high blood pressure, heart failure, or selected kidney and liver conditions but may disturb salts or cause dehydration.

\medterm{Doctor Of Osteopathic Medicine} A physician who has completed an accredited osteopathic medical degree and residency training and is licensed to diagnose, prevent, and treat disease.


\medterm{Drug Addiction (Substance Use Disorder)} A chronic disorder involving compulsive substance use despite harmful consequences, with impaired control, cravings, tolerance, or withdrawal.

\synonyms
Chemical Dependency
Substance Abuse Disorder

\medterm{Drug Administration Route} The pathway by which a medicine enters the body, such as oral, intravenous, intramuscular, inhaled, topical, or subcutaneous administration.

\medterm{Drug Allergies} Drug allergies are immune-mediated reactions to a medicine, ranging from rash to anaphylaxis; they must be distinguished from predictable side effects and nonimmune intolerance.

\medterm{Drug Allergy} An adverse immune-mediated reaction to a medication, representing a hypersensitivity
response that is distinct from side effects, toxicities, and idiosyncratic reactions.

\medterm{Drug Antagonism} An interaction where one drug blocks or decreases the effectiveness of another, such as
when an antidote blocks a poison.

\medterm{Drug Bioavailability} The proportion of a dose that reaches the bloodstream in an active form. It depends on how the medicine is given, how well it is absorbed, and how much is broken down before reaching the blood; intravenous doses enter the bloodstream completely.

\medterm{Drug Classification Systems} Systems that group medicines according to properties such as their chemical structure, mechanism of action, therapeutic use, or effects on the body.

\medterm{Drug Clearance} The body’s ability to remove a medicine from the bloodstream, mainly through the liver and kidneys. Clearance affects how long a medicine works and how often it should be taken; liver or kidney disease may require dose adjustment.

\medterm{Drug Dependence} A state of physiological or psychological adaptation to a drug, characterized by a
withdrawal syndrome upon discontinuation or dose reduction. Physical dependence: the
body has adapted to the presence of the drug, and abrupt cessation produces a
predictable withdrawal syndrome specific to the drug class.

\medterm{Drug Emulsion} A medicine made by dispersing tiny droplets of one liquid throughout another liquid that normally would not mix with it. An added stabilizer keeps the mixture uniform so the medicine can be given safely.

\medterm{Drug Eruption} A skin eruption caused by a medication, ranging from a mild morbilliform rash to severe reactions such as Stevens–Johnson syndrome or toxic epidermal necrolysis.

\medterm{Drug Hypersensitivity Reactions} Immune-mediated adverse reactions to medicines. They may be immediate or delayed and can range from a mild rash to severe systemic illness.

\medterm{Drug Hypersensitivity Syndrome} A serious delayed reaction to a medicine, usually beginning two to eight weeks after it is started. It may cause fever, widespread rash, facial swelling, enlarged lymph nodes, abnormal blood cells, and injury to the liver, kidneys, lungs, or heart.

\medterm{Drug Implants} Implantable dosage forms or devices that release a medicine over time; selection and monitoring depend on the product and clinical indication.

\medterm{Drug Induced Liver Disease} Drug-induced liver injury occurs when a medicine, supplement, or toxin damages liver cells or bile flow. It may cause fatigue, nausea, itching, dark urine, jaundice, or no symptoms; identifying and stopping the responsible exposure under medical guidance is central to care.

\medterm{Drug Interaction} A change in the effect or concentration of one medicine caused by another medicine, food, substance,
or disease.

\medterm{Drug Metabolism} The enzymatic conversion of medicines into metabolites, mainly in the liver but also in the intestine, kidneys, lungs, and other tissues.

\medterm{Drug Metabolism Pharmacogenomics} The study of how inherited genetic differences change the way a person breaks down medicines. This information can sometimes help clinicians choose a safer or more effective medicine or dose.

\medterm{Drug Overdose} Harm caused by taking too much medicine or taking it by the wrong route. Effects range from sleepiness, vomiting, and confusion to seizures, organ damage, or stopped breathing. Suspected overdose requires urgent poison-control or emergency advice.

\medterm{Drug Overdose Management} Management of poisoning begins with assessment and support of airway, breathing,
circulation, and neurologic status. Decontamination, enhanced elimination, and specific antidotes are considered only
when appropriate for the substance, timing, and patient.

\medterm{Drug Packaging} The container, labeling, and protective materials used to preserve a medicine and support correct identification, storage, and use.

\medterm{Drug Reaction With Eosinophilia And Systemic Symptoms} A severe hypersensitivity syndrome characterized by fever, rash, lymphadenopathy, and
multiorgan involvement triggered by medications, requiring prompt drug cessation and
systemic corticosteroids.

\medterm{Drug Resistance} The ability of microorganisms, cancer cells, or parasites to survive and proliferate
despite exposure to drugs that were previously effective.

\medterm{Drug Therapeutics} The use of medicines to prevent, diagnose, or treat disease while balancing expected benefit and
potential harm.

\medterm{Drug Tolerance} A decreased pharmacological response to a drug after repeated administration, requiring
dose escalation to achieve the same effect.

\medterm{Drug Toxicity} Harmful effects caused by a medicine or its metabolites. Toxicity may result from excessive exposure, accumulation, interactions, or an unpredictable individual reaction.

\medterm{Drug Use First Aid} Immediate care for suspected drug intoxication, including assessment of breathing and responsiveness, emergency activation, overdose-specific support, and naloxone when opioid toxicity is suspected.

\medterm{Drug Vehicle} An inactive substance or carrier used to deliver an active drug ingredient in a formulation.

\medterm{Drug Withdrawal} A clinically significant syndrome or return of symptoms caused by abrupt reduction or
cessation of a drug after physiologic adaptation or dependence has developed. Tapering
is required for some medications to reduce withdrawal and rebound effects.

\medterm{Drug, Prescription} A medicine legally supplied under the order of an authorized prescriber, with its indication, dose, duration, interactions, and risks specified for the patient.

\medterm{Drug, Sulfa} A sulfonamide drug; some are antibacterial medicines, while others are used as diuretics or for other indications, and sulfonamide reactions vary by drug and patient.

\medterm{Drug-Drug Interactions} Changes in the effect or amount of one medicine caused by another medicine. The combination may make treatment less effective or increase side effects, so all prescription, nonprescription, and supplement use should be reviewed.

\medterm{Drug-Eluting Balloon (Vascular)} A vascular balloon catheter that delivers a medicine to the vessel wall during dilation,
without leaving a permanent implant.

\synonyms
Vascular

\medterm{Drug-Eluting Stent} A vascular stent that slowly releases a medicine to reduce tissue growth and the risk of restenosis after angioplasty.

\medterm{Drug-Eluting Stent (Vascular)} A vascular stent coated with medicine that is released gradually to reduce tissue growth
and recurrent narrowing.

\synonyms
Vascular

\medterm{Drug-Facilitated Sexual Assault} Sexual contact occurring when a person's capacity to consent or resist is impaired by
alcohol, medication, or another drug, whether administered covertly or knowingly
consumed. Medical care includes urgent safety assessment, treatment of injuries,
toxicology sampling when indicated, prophylaxis, and trauma-informed support.

\medterm{Drug-Induced Autoimmunity} An immune reaction caused by a medicine in which the body attacks its own tissues. It may resemble lupus, cause joint pain, fever, rash, or organ inflammation, and often improves after the responsible medicine is stopped under medical guidance.

\medterm{Drug-Induced Diarrhea} Diarrhea caused by a medicine or drug interaction. Common causes include antibiotics, laxatives, metformin, and medicines that alter intestinal motility.

\medterm{Drug-Induced Immune Hemolytic Anemia} An uncommon immune-mediated destruction of red blood cells triggered by a medication; evaluation includes blood tests and review of exposures, with specialist management.

\medterm{Drug-Induced Liver Injury} Liver damage caused by a medicine, herbal product, or dietary supplement. It may produce
abnormal liver tests, hepatitis, cholestasis, or liver failure; evaluation requires considering
the exposure, timing, competing causes, and severity.

\medterm{Drug-Induced Low Blood Sugar} Low blood glucose caused by a medicine, most often insulin or an insulin secretagogue. Symptoms may include sweating, tremor, confusion, seizure, or loss of consciousness.

\medterm{Drug-Induced Lupus Erythematosus} A lupus-like autoimmune syndrome caused by certain medicines. Symptoms commonly include joint pain, muscle pain, and fever and usually improve after the causative medicine is stopped.

\medterm{Drug-Induced Pulmonary Disease} Lung injury caused by a medicine or treatment. It may produce cough, breathlessness, fever, or abnormal lung imaging; the medicine, timing, and other possible causes must be reviewed promptly because some cases can progress to respiratory failure.

\medterm{Drug-Induced Thrombocytopenia} A low platelet count caused by a medicine through immune destruction, reduced production, or other mechanisms, increasing bleeding risk.

\medterm{Drug-Induced Tremor} Tremor caused or worsened by a medicine, stimulant, or toxin. Assessment considers timing, dose, other causes, and whether the medicine can be changed.

\medterm{DTaP Vaccine} A vaccine against diphtheria, tetanus, and pertussis (whooping cough), given mainly to infants and young children in a multi-dose series. It reduces severe illness, complications, and spread; schedules vary by country.

\medterm{Eflornithine} A medicine that slows cell growth by blocking an enzyme needed to make certain cell components. A skin cream can reduce unwanted facial hair; an intravenous form has been used for selected African trypanosome infections.


\medterm{Enzyme Inhibitor} A substance that slows or blocks an enzyme, a protein that helps chemical reactions occur in the body. Medicines use this action to change processes such as blood pressure, inflammation, infection, hormone production, or cancer-cell growth.


\medterm{Estrogen Inhibitor} A medicine that reduces estrogen production or blocks estrogen’s effects. It may treat selected hormone-sensitive cancers or other conditions; side effects can include hot flashes, bone loss, joint symptoms, or blood clots.

\medterm{Etidronate} A bisphosphonate medicine that slows the breakdown and rebuilding of bone. It is used selectively for some disorders of abnormal bone remodeling; treatment depends on the condition and can affect the stomach, kidneys, or minerals.

\medterm{Evidence-Based Medicine} The use of reliable research together with clinical expertise and the patient’s values and circumstances to make healthcare decisions. It supports, but does not replace, discussion between patient and clinician.

\medterm{Expectorant} A medicine intended to loosen or thin mucus so it is easier to cough up. Benefit varies with the cause of cough and the product; persistent breathlessness, fever, chest pain, or blood needs medical assessment.

\medterm{Extensively Drug-Resistant TB} Tuberculosis caused by bacteria resistant to rifampicin and usually isoniazid, at least one fluoroquinolone, and at least one of the key drugs bedaquiline or linezolid. Diagnosis requires drug-susceptibility testing and specialist treatment.

\medterm{Family Medicine} A medical specialty providing comprehensive continuing care for people of all ages, including prevention, acute illness, chronic disease, and coordination of referrals.

\medterm{Fentanyl} A potent synthetic opioid analgesic used for severe pain and anesthesia. It can cause respiratory depression, sedation, constipation, nausea, tolerance, dependence, and fatal overdose, especially when combined with alcohol or other sedatives.

\medterm{Fertility Agent} A medicine or biologic intervention intended to improve the chance of conception by inducing ovulation, supporting sperm production, or treating a reproductive mechanism.

\medterm{Fidaxomicin} A narrow-spectrum oral macrocyclic antibiotic used to treat Clostridioides difficile infection. It has minimal systemic absorption and reduces disruption of normal intestinal flora compared with many alternatives, but recurrence can still occur.

\medterm{Filgrastim} A recombinant granulocyte colony-stimulating factor that increases production and release of neutrophils from bone marrow. It is used to reduce or treat selected forms of chemotherapy-related neutropenia; bone pain and rare splenic complications can occur.

\medterm{Finasteride} A medicine that lowers the conversion of testosterone to dihydrotestosterone. It can shrink an enlarged prostate and help some types of hair loss but may reduce libido, affect sexual function, or alter prostate-test results.

\medterm{Fixed Drug Eruption} A medicine-related skin reaction that returns in the same place whenever the person takes the responsible medicine. It may appear as a sharply outlined red or dark patch, sometimes with a blister or sore, and can leave lasting discoloration.

\medterm{Fluticasone} A corticosteroid available in inhaled, nasal, topical, and other formulations to reduce inflammation. Adverse effects depend on route and may include local irritation, oral candidiasis, skin thinning, adrenal suppression, or growth effects with prolonged exposure.

\medterm{Food-Drug Interaction} A change in how a medicine works or is absorbed because of food, drink, nutrients, or supplements. Examples include grapefruit affecting some medicines, vitamin K changing warfarin’s effect, and calcium reducing absorption of some antibiotics.

\medterm{Forensic Medicine} Application of medical knowledge to legal questions, including injury interpretation, cause of death, toxicology, identification, sexual assault evaluation, and documentation.

\medterm{Forensic Medicine (Legal Medicine)} Branch of medicine that applies medical, biological, and scientific knowledge to legal
questions and the administration of justice.

\synonyms
Legal Medicine

\medterm{Foscarnet} An intravenous antiviral agent used for selected cytomegalovirus or herpesvirus infections, especially when resistance or intolerance limits other treatment. Important toxicities include kidney injury and disturbances of calcium, magnesium, potassium, or phosphate.

\medterm{Furosemide} A diuretic that makes the kidneys remove more salt and water in urine. It treats swelling and some blood-pressure or heart conditions but may cause dehydration, abnormal body salts, kidney problems, or hearing damage.

\medterm{Ganglionic Blocker} A medicine that interrupts nerve signals at switching points controlling automatic body functions. It can affect blood pressure, heart rate, digestion, sweating, vision, and bladder function, so these medicines are rarely used clinically.

\medterm{Generic Drug} A medicine equivalent to a brand-name medicine in active ingredient, strength, dosage form, and intended use, usually sold under a nonproprietary name.

\synonyms
Generic medicine

\medterm{Generic Medicine} A medicine with the same active ingredient, strength, form, route, and intended effect as an approved reference medicine. It may look different or have different inactive ingredients but must meet regulatory quality and performance requirements.

\medterm{Generic Name, Drug} The official nonproprietary name of a medicine's active ingredient, used consistently across manufacturers and countries. It identifies the medicine itself rather than a company's brand name and helps prevent confusion when products have different trade names.

\medterm{Genomic Medicine} Use of a person’s genetic or genomic information to guide disease diagnosis, risk assessment, prevention, or treatment.

\medterm{Geriatric Medicine} The medical specialty focused on the health, function, diseases, medications, and goals of older adults, often using comprehensive and coordinated care.

\medterm{Growth Inhibitor} A medicine or biological substance that slows or stops cell growth. In cancer treatment, it may block cell division or a growth signal, but it can also affect healthy cells and cause side effects.

\medterm{Haemophilus Vaccine} A vaccine that protects against serious disease caused by Haemophilus influenzae type b, including meningitis, pneumonia, and bloodstream infection. It is given mainly in infancy according to the recommended immunization schedule.

\medterm{Hepatic Drug Metabolism} The liver changes medicines into forms that the body can use or remove. This may stop a medicine’s effect, activate a medicine that works after conversion, or produce a harmful product; age, illness, genetics, and other medicines can alter the process.

\medterm{Hepatitis B Vaccine} A recombinant vaccine that prevents hepatitis B infection, which can cause chronic liver disease, cirrhosis, and liver cancer. It is given in a series beginning at birth in many schedules; infants exposed at birth may also need hepatitis B immune globulin.

\medterm{Hepatitis B Vaccine (Birth Dose)} The first hepatitis B vaccine dose, given soon after birth, helps prevent infection passed from a parent to the newborn. If the birth parent has hepatitis B, the infant also needs hepatitis B immune globulin promptly; local schedules determine later doses.

\synonyms
Birth Dose

\medterm{HER2 Inhibitor} A targeted cancer treatment that blocks signaling through human epidermal growth factor receptor 2 in cancers with HER2 overexpression or amplification.

\medterm{High Blood Pressure - Medicine-Related} Hypertension caused or worsened by a medicine or substance, such as NSAIDs, steroids, decongestants, stimulants, oral contraceptives, or immunosuppressants; review and substitution may improve control.

\medterm{Histamine-2 Receptor Antagonist} A medicine that reduces stomach-acid production by blocking histamine signals in the stomach. It may treat reflux, indigestion, or ulcers; headache, diarrhea, confusion, and medicine interactions can occur.

\medterm{Hormone Inhibitor} A medicine or biologic that blocks hormone synthesis, release, receptor binding, or downstream signaling, used in selected endocrine disorders and hormone-sensitive cancers.

\medterm{HPV Vaccine} A vaccine that prevents infection with human papillomavirus types that cause several cancers and genital warts. It is routinely offered in early adolescence, can start from age 9, and may require two or three doses depending on age and immune status; it does not treat existing HPV infection.

\medterm{Hypoglycemic Agent} A medicine that lowers blood glucose, including insulin, insulin secretagogues, insulin sensitizers, incretin-based drugs, and sodium--glucose cotransporter inhibitors; choice depends on diabetes type and patient risks.

\medterm{Hypolipidemic Agent} A medicine that lowers fats in the blood, such as LDL cholesterol or triglycerides. It helps reduce cardiovascular risk when combined with appropriate diet, activity, and treatment of other risk factors.


\medterm{Immunosuppressant} A drug that reduces immune system activity, used to prevent organ transplant rejection
and to treat autoimmune and inflammatory diseases.

\medterm{Immunosuppressant Drug} A medicine that reduces immune-system activity, often to prevent transplant rejection or treat autoimmune disease.

\medterm{Immunosuppressive Drug} A drug that suppresses immune responses, used for organ transplantation,
graft-versus-host disease, and autoimmune/inflammatory diseases.

\medterm{Inactivated Polio Vaccine} An injected vaccine made from killed polioviruses. It trains the immune system to prevent poliomyelitis and is used in routine childhood immunization programs; the schedule and use of other polio vaccines vary by country.

\medterm{Influenza Vaccine} A seasonal vaccine that reduces the risk of influenza and its complications. It is updated as needed and is offered each year to most people from early infancy; some young children need two doses in their first vaccinated season, according to local recommendations.

\medterm{Inotrope} A drug that changes myocardial contractility; positive inotropes increase the force of contraction. Dobutamine and milrinone may be used short term under specialist monitoring for cardiogenic shock or selected cases of acute decompensated heart failure with low output. Dopamine, epinephrine, and norepinephrine are primarily vasopressors or emergency agents with additional inotropic effects, while digoxin is a separate cardiac glycoside used in selected chronic heart-failure or rate-control settings. These medicines can cause arrhythmias, ischemia, or other serious effects and are not routine long-term treatment for chronic heart failure.


\medterm{Intensive Care Medicine} Specialty managing critically ill patients. Components: monitoring, ventilatory support,
hemodynamic support, infection control, nutrition, sedation, pain management. Common
conditions: sepsis, ARDS, traumatic brain injury, post-cardiac arrest, multiorgan
failure.

\medterm{Internal Medicine} The medical specialty focused on prevention, diagnosis, and nonsurgical treatment of diseases in adults, including complex multisystem illness.

\medterm{Japanese Encephalitis Vaccine} A vaccine that protects against Japanese encephalitis, a mosquito-borne viral infection of
the central nervous system. Vaccine type, schedule, and eligibility depend on age,
destination or exposure risk, and local recommendations.

\medterm{Laxative} A medicine that helps empty the bowel and is used for constipation or bowel preparation. Some add fiber, some draw water into stool, and others stimulate bowel movement; the best choice depends on the cause, other medicines, and the person’s health.

\medterm{Legal Medicine} The application of medical knowledge to legal questions, including forensic examination, injury assessment, death investigation, and issues of capacity or consent.


\medterm{Liquid Medication Administration} Giving a medicine in liquid form by mouth, feeding tube, injection, or another prescribed route. The dose must be measured accurately with the supplied device because liquid strengths differ; household spoons can cause dosing errors.

\medterm{Local Anesthetic} A medicine that temporarily numbs a specific area without usually causing loss of consciousness. It blocks pain signals in nearby nerves and is used for procedures, dental work, wound repair, and some forms of regional pain relief.

\medterm{Luteinizing Hormone-Releasing Hormone Agonist} A medicine that first stimulates and then, with continued use, suppresses signals from the pituitary gland that control the ovaries or testes. It is used for selected cancers, endometriosis, fibroids, and puberty disorders.

\synonyms
GnRH agonist; LHRH agonist

\medterm{Macrolide Antibiotic} An antibiotic of the macrolide class that binds the bacterial 50S ribosomal subunit and usually inhibits protein synthesis; activity and resistance vary by organism and drug.

\medterm{Malaria Vaccine} A vaccine intended to reduce malaria caused by \textit{Plasmodium} parasites, especially \textit{P. falciparum}; the RTS,S/AS01 and R21/Matrix-M vaccines are used in recommended childhood immunization programs in some endemic settings, alongside other preventive measures.


\medterm{MAO Inhibitor} A medicine that slows breakdown of certain brain chemicals and is used for selected depression or Parkinson disease. It can interact dangerously with other medicines and some foods, so it requires careful prescribing and instructions.


\medterm{Measles Vaccine} Measles vaccine is a live attenuated vaccine that induces active immunity against measles virus. It is usually given as part of the measles--mumps--rubella (MMR) or measles--mumps--rubella--varicella (MMRV) vaccine.

\medterm{Measles-Rubella Vaccine} A combined live attenuated vaccine containing measles and rubella antigens without the
mumps component, widely used in countries where mumps is not a significant public health
burden, including India's Universal Immunization Schedule.

\medterm{Medication} A substance or biological product used to prevent, diagnose, treat, or control a disease or symptom. Medicines may be prescribed, bought without a prescription, applied to the body, inhaled, swallowed, or given by injection.

\medterm{Medication Adherence} The extent to which a patient’s medication-taking behavior corresponds with the agreed
treatment plan. Adherence is influenced by understanding, adverse effects, cost, access,
regimen complexity, beliefs, and cognitive or social factors.

\medterm{Medication Adverse-Effect Profile} The clinically relevant pattern of predictable pharmacologic effects and harmful
reactions associated with a medication. It includes common dose-related effects, serious
rare events, organ-specific toxicity, interactions, and monitoring requirements.

\medterm{Medication Cascade} A prescribing sequence in which an adverse drug effect is misinterpreted as a new
medical condition and another drug is prescribed to treat it. Recognizing cascades is an
important goal of medication reconciliation and deprescribing.

\medterm{Medication Deprescribing} A planned, supervised process of reducing the dose or stopping a medicine when its
current or future harms outweigh its benefits or when treatment goals change. Tapering
may be required to prevent withdrawal or rebound disease.

\medterm{Medication Error} A preventable event that may cause or lead to inappropriate medication use or patient
harm while the medication is in the control of a healthcare professional, patient, or
consumer. Errors can occur during prescribing, transcribing, dispensing, administration,
or monitoring.

\medterm{Medication Management} The systematic process of optimizing pharmacotherapy to achieve therapeutic goals while
minimizing adverse effects, encompassing prescribing, dispensing, administration,
monitoring, and deprescribing.

\medterm{Medication Overuse Headaches} Frequent headaches caused or worsened by regular overuse of headache medicines.

\synonyms
Rebound Headache

\medterm{Medication Reconciliation} A structured process of obtaining and comparing the most accurate medication list with
current prescriptions at transitions of care. It identifies omissions, duplications,
incorrect doses, interactions, and unintended changes.

\medterm{Medication Review} Systematic examination of all medications for appropriateness. Components: indication,
efficacy, safety, adherence, interactions. Tools: Beers Criteria, STOPP/START. Goals:
deprescribing, optimizing therapy, reducing adverse effects.

\medterm{Medication Safety} Practices that reduce medication errors and prevent avoidable harm during prescribing, dispensing, administration, and use.

\medterm{Medication Side Effects} Unintended harmful responses to medications occurring at standard doses, ranging from
mild (nausea, rash, dry mouth) to severe (anaphylaxis, Stevens-Johnson syndrome, QT
prolongation, hepatotoxicity).

\medterm{Meningococcal Vaccine} Conjugate vaccine (MenACWY) and serogroup B vaccine (MenB) given at 11-12 years with
booster at 16 years. Protects against Neisseria meningitidis, causing meningitis and
septicemia. MenACWY booster at 16 years covers the high-risk period for serogroup C and
Y disease.

\medterm{Mineralocorticoid Receptor Antagonist} A medicine that blocks aldosterone, a hormone that makes the body retain salt and water. It is used for selected heart failure, high blood pressure, and hormone disorders but can raise blood potassium and worsen kidney problems.

\medterm{MMR Vaccine} A live attenuated vaccine against measles, mumps, and rubella. It is given according to the
immunization schedule and is generally avoided during pregnancy and in people with certain severe immune deficiencies.

\medterm{Molecular Medicine} Molecular medicine applies molecular and cellular biology to understand, diagnose, prevent, and treat disease. It includes genomics, precision oncology, molecular biomarkers, targeted therapy, and other mechanism-based approaches.

\medterm{Monoamine Oxidase Inhibitor} A medicine that blocks monoamine oxidase, an enzyme that breaks down certain brain chemicals. It can treat selected depression or neurologic disorders but has important interactions with other medicines and some foods.

\medterm{MRNA Vaccine} A vaccine that gives cells temporary instructions to make a harmless piece of a germ. The immune system learns to recognize that piece and respond later; the messenger RNA does not enter the cell nucleus or change DNA.

\medterm{Multi-Drug Resistance} Multidrug resistance is resistance of a microorganism, cancer, or other disease process to multiple drugs that would otherwise be active. Mechanisms include target changes, drug inactivation, reduced uptake, efflux, and altered repair or survival pathways.

\medterm{Multidrug-Resistant Extensively Drug-Resistant TB} Multidrug-resistant tuberculosis is resistant to key first-line drugs, while extensively drug-resistant TB has additional resistance that greatly limits options. Diagnosis requires molecular or culture testing, and treatment uses prolonged individualized combinations with public-health oversight.

\medterm{Neuropathy Secondary To Drugs} Drug-induced neuropathy is peripheral nerve damage caused by a medicine or toxin, often producing numbness, burning pain, weakness, or altered reflexes; management may require dose reduction, withdrawal, and treatment of symptoms.

\medterm{Nonsteroidal Anti-Inflammatory Drug} A medicine that inhibits cyclooxygenase enzymes and reduces prostaglandin production, producing analgesic,
anti-inflammatory, and fever-reducing effects. NSAIDs can cause gastrointestinal,
kidney, cardiovascular, or other adverse effects.

\medterm{NSAIDs} Nonsteroidal anti-inflammatory drugs are a class of drugs that inhibit cyclooxygenase
enzymes, COX-1 and COX-2, reducing the synthesis of prostaglandins and thromboxanes from
arachidonic acid. They produce analgesic, anti-inflammatory, antipyretic, and
antiplatelet effects.


\medterm{Occupational Medicine} The medical specialty concerned with the effects of work on health, prevention of work-related illness and injury, and safe return
to work.

\medterm{Opioid Analgesics} Strong pain medicines that act on opioid receptors in the brain, spinal cord, and other tissues. They can relieve severe pain but may cause sleepiness, constipation, dependence, and dangerously slow breathing or overdose.

\medterm{Opioid Antagonist} A medicine that blocks opioid receptors and can reverse opioid effects such as respiratory depression.
Naloxone is a commonly used example.

\medterm{Oral Hypoglycemic} A medicine taken by mouth to lower blood glucose, used mainly in type 2 diabetes. Different medicines improve insulin sensitivity, increase insulin release, slow glucose absorption, or increase glucose loss in urine; choice depends on the patient’s risks and kidney function.

\medterm{Oral Medicine} Specialty diagnosing oral mucosal diseases, dental causes of facial pain, oral
manifestations of systemic disease.

\medterm{Oral Polio Vaccine} A live attenuated poliovirus vaccine administered by mouth. Its formulation and use
depend on the public-health program and circulating poliovirus strains; it has played an important role in polio control.


\medterm{Orphan Drug} A pharmaceutical agent developed specifically to treat a rare medical condition
(affecting fewer than 200,000 people in the United States, as defined by the Orphan Drug
Act of 1983).

\medterm{Over-The-Counter (OTC) Drug} A medicine sold without a prescription for uses and doses described on its label. “Over the counter” does not mean risk-free; incorrect use, interactions, or excessive doses can still cause serious harm.

\medterm{Over-The-Counter Drug} A medication that can be purchased without a prescription, deemed safe and effective for
self-medication when used according to label instructions. OTC drugs have established
safety profiles, low toxicity, and clear dosing guidelines suitable for consumer use.

\medterm{P2Y12 Inhibitor} An antiplatelet drug that blocks the platelet P2Y12 ADP receptor, reducing activation
and aggregation. Clopidogrel, prasugrel, and ticagrelor are used in selected acute
coronary syndrome and percutaneous coronary intervention settings, with bleeding and
drug-specific contraindications considered.

\medterm{Palifermin} A laboratory-made form of a natural growth factor that helps mouth and throat lining cells repair themselves. It is used in selected people receiving intensive treatment before blood-forming stem-cell transplantation to reduce severe mouth sores.

\medterm{Palonosetron} A long-acting anti-nausea medicine that blocks serotonin signals involved in vomiting. It is used to prevent nausea and vomiting from chemotherapy or surgery; headache, constipation, and rare heart-rhythm changes can occur.

\medterm{Parasympatholytic} A medicine that blocks some signals of the parasympathetic nervous system. It may widen the airways, relax the bladder, reduce motion sickness, or speed a slow heartbeat, but can cause dry mouth, blurred vision, constipation, urinary retention, or confusion.

\medterm{Parasympathomimetic} A drug that mimics or enhances parasympathetic (cholinergic) activity. Direct agonists:
muscarinic (pilocarpine—glaucoma, dry mouth) and nicotinic agonists.

\medterm{Paroxetine} Paroxetine is a selective serotonin-reuptake inhibitor used for depression, anxiety disorders, obsessive-compulsive disorder, PTSD, and some other conditions; sexual dysfunction, withdrawal symptoms, serotonin syndrome, and pregnancy-related risks require consideration.

\medterm{Pediatric Nuclear Medicine} Special considerations for children: lower doses, sedation often needed, distraction
techniques. Dose reduction: weight-based activity, optimized protocols.

\medterm{Pefloxacin} A fluoroquinolone antibiotic that blocks bacterial enzymes needed to copy and repair DNA. Its use is limited by resistance and potentially serious effects involving tendons, nerves, the heart, and the nervous system.

\medterm{PEG} An abbreviation that may refer to polyethylene glycol, a laxative that draws water into stool, or to a polyethylene-glycol chemical group used to modify medicines. The exact meaning depends on the clinical context.


\medterm{Pegaspargase} A long-acting enzyme medicine used in combination treatment for some acute lymphoblastic leukemias. It lowers asparagine needed by leukemia cells; allergic reactions, pancreatitis, liver injury, clotting problems, and bleeding are important risks.

\medterm{Pegfilgrastim} A long-acting growth factor that stimulates the bone marrow to make neutrophils. It helps prevent infection caused by chemotherapy-related low white-cell counts; bone pain and rare spleen rupture or serious inflammation can occur.

\medterm{Penicillamine} A medicine that binds certain metals and can also reduce immune activity. It is used selectively for copper overload and some autoimmune or connective-tissue disorders; serious blood, kidney, skin, or allergic reactions can occur.

\medterm{Pentavalent Vaccine} A WHO-prequalified five-in-one combination vaccine containing Diphtheria toxoid, Tetanus
toxoid, whole-cell Pertussis antigen, Hepatitis B (HBsAg), and Haemophilus influenzae
type b (Hib) conjugate, administered as a single intramuscular injection.

\medterm{Pentobarbital} A barbiturate medicine that depresses brain activity and can cause sedation or control seizures. It has a high risk of slowed breathing, coma, dependence, and fatal overdose, especially with alcohol or other sedatives.


\medterm{Perioperative Antibiotic Prophylaxis} Administration of an appropriate antimicrobial within a defined interval before incision
to reduce infection risk for selected procedures. Choice, timing, redosing, and
discontinuation depend on procedure, organisms, patient factors, and local guidance.

\medterm{Pertussis Vaccine} A vaccine containing selected antigens from Bordetella pertussis, given in combination vaccines to prevent whooping cough and reduce severe disease and transmission.

\medterm{Pharmaceutical} Relating to medicines, including their development, preparation, manufacture, testing, storage, or use.

\medterm{Pharmacodynamic Interaction} A drug interaction in which one drug changes the effect of another at its site of action
without necessarily changing its concentration. Effects may be additive, synergistic, or
antagonistic, such as combined respiratory depression or bleeding risk.

\medterm{Pharmacokinetic Interaction} A drug interaction that changes absorption, distribution, metabolism, or excretion of
another drug, thereby altering its concentration over time. Enzyme inhibition, enzyme
induction, transporter effects, and altered gastric pH are common mechanisms.

\medterm{Phosphodiesterase Inhibitor} A medicine that blocks enzymes which normally break down chemical signals inside cells. Different types relax airways or blood vessels, improve heart pumping, or treat erectile dysfunction; effects and risks depend on the specific type.


\medterm{Pneumococcal Conjugate Vaccine} A vaccine that joins parts of the pneumococcus bacterium to a carrier protein so the immune system can build lasting protection. It helps prevent pneumonia, meningitis, bloodstream infection, and other disease; product and schedule depend on age and risk.

\medterm{Pneumococcal Vaccine} A vaccine that protects against disease caused by Streptococcus pneumoniae, including invasive infection, pneumonia, meningitis, and otitis media; product and schedule depend on age and risk.

\medterm{Poliovirus Vaccine} A vaccine that trains the immune system to prevent poliovirus infection and paralysis. Inactivated injectable and live oral forms are used in different programs; the schedule and product depend on age, health, and local public-health guidance.

\medterm{Prazosin} An alpha-1 blocker that relaxes blood vessels and may relax the bladder outlet. It is used for high blood pressure and, in selected patients, trauma-related nightmares or urinary symptoms; dizziness and fainting are most likely after starting or increasing the dose.

\medterm{Pre-Medication} Medicine given before a procedure, treatment, or exposure to prevent a predictable problem such as pain, nausea, allergy, infection, or anxiety.


\medterm{Prescription Drug} A medication whose dispensing or use is subject to prescription requirements under the applicable law. It
requires professional oversight because of its risks, complexity, or need for monitoring;
regulatory approval and prescribing rules vary by country and product.

\medterm{Procalcitonin-Guided Antibiotic Therapy} Use of blood procalcitonin measurements, together with examination and other tests, to help decide whether to start or stop antibiotics in selected infections. The result cannot by itself confirm bacterial infection or safely replace clinical judgment.

\medterm{Prostaglandin Antagonist} A drug that blocks a prostaglandin receptor or pathway, reducing selected prostaglandin-mediated effects such as inflammation, pain, uterine contraction, or gastric secretion.

\medterm{Protease Inhibitor} An antiviral medicine that blocks a virus enzyme needed to cut newly made proteins into working parts. In HIV treatment, it prevents formation of mature infectious virus and is used with other medicines.

\medterm{Proton Pump Inhibitor} A medicine that strongly reduces stomach-acid production by blocking the acid pump in stomach cells. It treats reflux and ulcers; long-term or unnecessary use can cause diarrhea, infections, low magnesium, or vitamin deficiencies.

\medterm{Psychotropic Drug} A medicine that affects brain function and may change mood, thinking, behavior, sleep, or perception. The term includes antidepressants, antipsychotics, mood stabilizers, stimulants, and some medicines for anxiety or sleep.

\medterm{Psychotropic Medication} A medicine that changes brain activity and may affect mood, thoughts, behaviour, sleep, or perception. Examples include antidepressants, antipsychotics, mood stabilizers, and some medicines for anxiety or attention disorders.


\medterm{Rabies Vaccine} A vaccine containing inactivated rabies virus used for preexposure protection or postexposure prophylaxis after a possible rabies exposure, together with wound care and immune globulin when indicated.

\medterm{Rehabilitation Medicine} The medical specialty that helps people restore or adapt movement, communication, thinking, work, and daily activities after illness or injury. It is also called physiatry.

\synonyms
physiatry.

\medterm{Renal Drug Excretion} Removal of medicines from the body by the kidneys into urine. Kidney disease can slow this process, allowing some medicines to build up and requiring dose changes or closer monitoring.

\medterm{Renin Inhibitor} A blood-pressure medicine that blocks renin, an enzyme involved in narrowing blood vessels and retaining salt. It lowers blood pressure but can raise potassium, affect kidney function, and harm a developing fetus.

\medterm{Resistance Antibiotic (Antibiotic Resistance)} Alternate index label for Antibiotic Resistance.

\medterm{Rock (Drug)} A street term commonly used for crack cocaine, a smokeable form of cocaine. It produces intense, short-lived stimulation but can cause dependence, dangerous heart rhythms, heart attack, stroke, seizures, overheating, or severe agitation.


\medterm{Rotavirus Vaccine} An oral live attenuated vaccine that protects infants and young children against severe
rotavirus gastroenteritis. The product, schedule, and age limits depend on the immunization program; mild fever or
diarrhea may occur.

\medterm{Scabicide} A medicine applied to the skin or taken by mouth to kill scabies mites. Treatment usually includes all close household or sexual contacts and washing or isolating recently used clothing and bedding to prevent reinfection.

\medterm{Selective Serotonin Reuptake Inhibitor} A medicine that increases serotonin signalling between nerve cells by slowing serotonin reuptake. It treats depression and several anxiety-related disorders; nausea, sexual difficulties, withdrawal symptoms, and rarely serotonin toxicity can occur.

\medterm{SGLT2 Inhibitor} A medicine that makes the kidneys pass more glucose and salt into the urine. It lowers blood sugar and can protect the heart and kidneys in selected people, but may cause dehydration, genital infections, or rarely ketoacidosis.


\medterm{Skeletal Muscle Relaxant} A medicine that reduces excessive contraction of skeletal muscles. It may relieve short-term muscle spasm or longer-term spasticity, but sleepiness, weakness, dizziness, and impaired breathing can occur with some drugs.

\medterm{Sleep Drug} A medicine used to help a person fall asleep, stay asleep, or maintain a regular sleep pattern. Benefits and risks vary widely; drowsiness, falls, impaired thinking, dependence, and interactions with alcohol or sedatives are important concerns.


\medterm{Steroid} A group of chemicals with a characteristic structure. In medicine, steroids include corticosteroids that reduce inflammation and natural or replacement hormones such as cortisol, estrogen, progesterone, and testosterone.

\medterm{Steroid Abuse} Steroid misuse is non-prescribed or excessive use of corticosteroids or anabolic-androgenic steroids, which can cause endocrine suppression, psychiatric effects, infection, metabolic disease, liver injury, cardiovascular harm, or tissue damage.

\medterm{Steroid Hormones} Lipid-soluble hormones derived from cholesterol, including cortisol, aldosterone, estrogen, progesterone, and testosterone.

\medterm{Steroid Injections - Tendon, Bursa, Joint} Injection of a corticosteroid, often with local anesthetic, into or around a joint, bursa, or tendon sheath to reduce inflammation and pain; injections directly into tendons are generally avoided because of rupture risk.

\medterm{Tamsulosin} Tamsulosin is a selective alpha-1A blocker that relaxes prostate and bladder-neck smooth muscle to improve urinary symptoms of benign prostatic hyperplasia; dizziness, postural hypotension, and ejaculatory dysfunction may occur.

\medterm{TCAs} Tricyclic antidepressants are medicines that increase the effects of certain brain chemicals. They can treat depression, nerve pain, and some other conditions but may cause sleepiness, dry mouth, constipation, low blood pressure, abnormal heart rhythms, or overdose toxicity.

\medterm{Tdap Vaccine} Tetanus, Diphtheria, and acellular Pertussis booster given at 11-12 years and during
each pregnancy (27-36 weeks gestation). Replaces one dose of Td. Important for pertussis
protection, particularly for protecting newborns through maternal antibody transfer.
Also given to adolescents and adults to maintain tetanus and diphtheria immunity.


\medterm{Terazosin} An alpha-blocking medicine that relaxes the prostate and blood vessels. It can improve urine flow in prostate enlargement and lower blood pressure, but dizziness or fainting may occur, especially after the first dose or dose increases.

\medterm{Therapeutic Drug Levels} The range of medicine concentrations in blood associated with benefit while limiting toxicity for a particular drug. A measured level must be interpreted with the dose, timing, symptoms, kidney or liver function, and the person’s response.

\medterm{Therapeutic Drug Monitoring} Measuring selected medicine levels in blood to adjust the dose for the best balance of benefit and safety. It is most useful when small changes in level can cause treatment failure or serious toxicity.

\medterm{Thrombolytic} A medicine that helps dissolve a blood clot by changing plasminogen into plasmin, which breaks down fibrin. It may be used urgently for selected strokes, heart attacks, or other life-threatening clots when benefits outweigh bleeding risk.

\medterm{Thrombolytic Drugs For Heart Attack} Medicines that dissolve a clot blocking a coronary artery during selected heart attacks. They can restore blood flow when urgent artery-opening procedures are unavailable or delayed, but may cause dangerous bleeding and are not suitable for everyone.


\medterm{Transfusion Medicine} The medical specialty that selects, prepares, tests, and gives blood or blood components such as red cells, platelets, plasma, and cryoprecipitate. It includes matching, preventing transfusion reactions, and monitoring patients.

\medterm{Travel Medicine} Medical care concerned with prevention and treatment of health problems related to travel.

\medterm{Tricyclic Antidepressant Overdose} Poisoning with a tricyclic antidepressant can cause drowsiness, confusion, dry mouth, seizures, low blood pressure, and dangerous heart rhythms. It is an emergency requiring immediate poison-control or hospital treatment.


\medterm{Tumor Necrosis Factor Inhibitor} A biological medicine that blocks tumor-necrosis factor, an immune signal involved in inflammation. It treats selected arthritis, psoriasis, and bowel diseases but can increase serious infection risk, so screening and monitoring are important.


\medterm{Tyrosine Kinase Inhibitor} A targeted medicine that blocks enzymes carrying growth signals inside cells. It can slow cancers driven by abnormal signaling, but side effects vary and may include diarrhea, skin changes, liver injury, high blood pressure, or heart problems.

\medterm{Uricosuric Agent} A medicine that increases renal excretion of uric acid, lowering serum urate in selected patients with gout. Adequate kidney function and hydration are important, and urinary uric-acid stones may be a concern.

\medterm{Urine Drug Screen} A preliminary urine test that looks for selected medicines or recreational drugs and their breakdown products. Results can be falsely positive or negative and may need confirmation with a more specific laboratory test.

\medterm{Uterotonic} A medicine that makes the uterus contract. It is used under medical supervision to start or strengthen labor, prevent or treat heavy bleeding after birth, or complete selected pregnancy-related procedures; excessive contraction can harm the mother or fetus.

\medterm{Vaccine Adjuvant} A substance added to some vaccines to enhance and prolong the immune response to the antigen. Examples include aluminum salts and newer adjuvant systems.



\medterm{Vaccine Therapy} Treatment that uses a vaccine to stimulate an immune response against a disease. It may prevent infection or, in selected cancer treatments, help the immune system recognize abnormal cells; its benefits depend on the specific vaccine and condition.

\medterm{Varicella Vaccine} A live attenuated vaccine that prevents or reduces the severity of chickenpox. The number and timing of doses
depend on age, vaccination history, immune status, pregnancy, and local recommendations.
It is generally avoided during pregnancy and in people with certain immune deficiencies.



\medterm{Xanthine Oxidase Inhibitor} A medicine that lowers uric-acid production by blocking an enzyme involved in breaking down substances from cells and food. It is used mainly to prevent gout attacks and, in selected cases, complications from rapid cancer-cell breakdown.

\section*{Brand Medicines}
